The variability of events that occur in a game, together with the heterogeneity of goal-setting occasions that players and teams face at different moments of the game, requires a systematic approach to process event data and untangle team-style models. The proposed methodology for transforming event data into process-discovery models is summarized in Fig.~\ref{fig:team_traces_methodology}.

An \emph{event} $e=(a,\rho,\pi,x,y,t)$ records an on-ball action: its type $a$ (e.g.\ pass or carry), the acting player $\rho$, the team $\pi$, the pitch location $(x,y)$, and the time $t$. A \emph{match} is the time-ordered sequence of its events. A \emph{possession} is a maximal run of consecutive events by the same team. It starts when the team gains the ball and ends when the team loses it, the match stops, or a shot is taken. Each possession is one process instance, identified by a \emph{case id}.

\begin{figure}[t!]
\includegraphics[width=\textwidth]{fig1.pdf}
\caption{Process Discovery methodology. Team purpose traces are fitted into the discovery algorithm where team behavior patterns are described employing logical artifacts.}\label{fig:team_traces_methodology}
\end{figure}

A \emph{purpose filter} $\Phi$ is a predicate on possessions. It keeps only the possessions relevant to a chosen tactical question, for example possessions that start in the own half and reach the opponent box. Since the ground truth of a team's intention is not directly observable, $\Phi$ operationalizes purpose through observable starting conditions, events, locations, durations, or outcomes. Examples include possessions that reach a specific zone, possessions that start from the goalkeeper and end in a turnover in the team's own half, or possessions that start with a recovery in the team's own penalty box and reach the offensive third within 15 seconds. This filtering reduces variability before discovery and makes the resulting models easier to interpret.

To compare spatial behavior across data providers, every event is mapped to a standardized $105\times68$ pitch and assigned one of $|Z|=24$ zones through a function $\zeta$. The \emph{activity} of an event is the pair $\alpha(e)=(a(e),\zeta(e))$, so each process step represents an action type performed in a field zone. Spatial aggregation is needed because raw event coordinates can be noisy and can make similar possessions appear different. A simple grid partition allows the discovery algorithm to identify larger-scale similarities in field use, as illustrated in Fig.~\ref{fig:purpose_traces_example}. Other partitions could be used depending on the analytical purpose, including coach-defined channels, tactical zones, hierarchical partitions, or trace-clustering approaches \cite{Clijmans2022,McCarthy2023}; here, the grid offers a transparent and computationally simple default.

\begin{figure}[t!]
\centering
\includegraphics[width=0.5\textwidth]{fig2.pdf}
\caption{Similar team traces starting on the defensive side of the field and being able to reach the opponent's penalty box. While the traces look similar, the raw location of the events might induce misleading differences. A simple field partitioning helps highlight the common zones the team used to build the sequence up and access the opponent penalty box.}\label{fig:purpose_traces_example}
\end{figure}

A possession then becomes a \emph{trace} $\sigma=\langle\alpha(e_1),\dots,\alpha(e_n)\rangle$. The multiset of all traces kept by $\Phi$ for a team is its \emph{event log} $L$. Process discovery learns a model $M=\mathcal{D}(L)$ from the log; we use the Heuristic Miner for $\mathcal{D}$ (Section~\ref{subsec:processdiscovery}). \emph{Conformance} then scores the pair $(M,L)$ on fitness, precision, generalization, and simplicity (Section~\ref{subsec:evaluation}). The task of this paper is, for a chosen purpose $\Phi$, to discover $M$ and read it as a description of how a team executes that purpose.
